\begin{otherlanguage}{arabic}
\large

\chapter*{%
  \begin{flushright}
    ملخص
  \end{flushright}%
}

مع تحوّل الأنظمة الكهروضوئية \textenglish{(PV)} إلى عنصر أساسي في أنظمة الطاقة الحديثة، تغدو الموثوقية المستدامة في ظل ظروف التشغيل الميدانية أمراً جوهرياً لقيمتها الاقتصادية. فالمحطات الكهروضوئية معرّضة لأعطال قد لا تكون مرئية بالفحص المباشر، ولا تظهر إلا من خلال انحرافات في القياسات الكهربائية والجوية. ومن ثَمّ يصبح الكشف والتشخيص المبكّران ضروريين للحفاظ على الأداء، والحدّ من التدهور، ودعم المراقبة على نطاق واسع.

تطوّر هذه الأطروحة نظاماً ناجعا للكشف عن الأعطال وتشخيصها في المحطات الكهروضوئية اعتماداً على قياسات ميدانية حقيقية، قابلاً للنشر على عتاد طرفي مدمج. ويتّبع النظام بنية من مرحلتين: تحدّد مرحلة كشف الإختلال الانحرافات عن السلوك الطبيعي، فيما تحدّد مرحلة تصنيف الأعطال نوع العطل المقابل. وتدعم واجهة مراقبة آنية المتابعة العملية.

وتتمثّل المساهمة المنهجية الرئيسية في \textenglish{PC-Flow}، وهو نموذج تدفّق تطبيعي باستدلال فيزيائي، مصمَّمة لكشف الإختلال بدقّة وكفاءة عالية. فبأقل من \textenglish{25,000} معلمة، يحقّق \textenglish{PC-Flow} قيمة \textenglish{PR-AUC} ثنائية مقدارها \textenglish{0.972}، وينجز الاستدلال خلال \textenglish{970}~ميكروثانية عند الملين التاسع والتسعين. أمّا في تصنيف الأعطال، فيبلغ مصنّف \textenglish{ExtraTrees} المقترَح قيمة \textenglish{F1}-ماكرو قدرها \textenglish{0.959}، ودقّة متوازنة تبلغ \textenglish{94.3\%}. وعلى المنصّة المدمجة، تبقى سلسلة المعالجة الكاملة دون \textenglish{31}~مللي ثانية عند الملين التاسع والتسعين.

وتُثبت هذه النتائج أنّ تحقيق إشراف دقيق ومنخفض الكمون على الأعطال ممكنٌ مباشرةً في موقع المحطة، بما يقلّص الفجوة الزمنية بين وقوع العطل واستجابة المشغّل، ويحدّ من خسائر الإنتاج المتراكمة، ويجعل المراقبة المستمرّة للأنظمة الكهروضوئية قابلة للتطبيق في البيئات الميدانية التي تشتدّ فيها الحاجة إليها.

\vspace{1cm}

\noindent\rule[2pt]{\textwidth}{0.5pt}
\textbf{الكلمات المفتاحية:} الأنظمة الكهروضوئية، الكشف عن الأعطال وتشخيصها، التدفّقات التطبيعية، النماذج المسترشدة بالفيزياء، الحوسبة الطرفية.

\noindent\rule[2pt]{\textwidth}{0.5pt}

\end{otherlanguage}
